% Part: incompleteness
% Chapter: theories-computability
% Section: introduction

\documentclass[../../../include/open-logic-section]{subfiles}

\begin{document}

% SEGMENT OLP-0302-S01

\olfileid{inc}{tcp}{int}

\olsection{அறிமுகம்}

\begin{editorial}
இந்தப் பிரிவு மீண்டும் எழுதப்பட வேண்டும்.
\end{editorial}

பின்வருவன நம்மிடம் உள்ளன:
\begin{enumerate}
\item ஒரு சார்பு~$\Th{Q}$ இல் பிரதிநிதித்துவப்படுத்தத்தக்கது என்பதன்
  வரையறை (\olref[req][int]{defn:representable-fn})
\item ஓர் உறவு~$\Th{Q}$ இல் பிரதிநிதித்துவப்படுத்தத்தக்கது என்பதன்
  வரையறை (\olref[req][rel]{defn:representing-relations})
\item $\Th{Q}$ இல் பிரதிநிதித்துவப்படுத்தத்தக்க சார்புகள் துல்லியமாகக்
  கணிக்கத்தக்க சார்புகளே என்று கூறும் ஒரு தேற்றம்
  (\olref[req][int]{thm:representable-iff-comp})
\item $\Th{Q}$ இல் பிரதிநிதித்துவப்படுத்தத்தக்க உறவுகள் துல்லியமாகக்
  கணிக்கத்தக்க உறவுகளே என்று கூறும் ஒரு தேற்றம்
  % SOURCE CORRECTION TA-INC-030: balance the fourth list item's reference parenthesis.
  (\olref[req][rel]{thm:representing-rels})
\end{enumerate}

ஒரு {\em கோட்பாடு} என்பது வருவித்தலின்கீழ் மூடிய வாக்கியங்களின் ஒரு
கணம்; அதாவது, $T$ ஆனது~$!A$ ஐ நிறுவும்போதெல்லாம் $!A$ ஆனது~$T$ இல்
இருக்கும் பண்பைக் கொண்ட கணம். ஒரு கோட்பாட்டை, சில வாக்கியங்களும்
அவற்றால் உட்கிடையாகப் பெறப்படும் அனைத்தும் சேர்ந்த ஒரு தொகுப்பாகக்
கருதுவது சிறந்தது. இனி, \olref[req][int]{sec} இல் உள்ள எட்டு
அடிகோள்களிலிருந்து வருவிக்கத்தக்க வாக்கியங்களின் கணமாகிய
{\em கோட்பாட்டையே}~$\Th{Q}$ என்று குறிப்பிடுவோம். $\Th{Q}$ இன்
வாய்பாடுகளை எண்களாகக் குறியீடாக்க முடியும் என்பதை நினைவுகூர்க;
$!A$ அத்தகைய ஒரு வாய்பாடு எனில், $!A$ ஐக் குறியீடாக்கும் எண்ணை
$\Gn{!A}$ குறிக்கட்டும். இக்குறியீடாக்கத்தைப் பொறுத்து, வாய்பாடுகளின்
பல கணங்கள் கணிக்கத்தக்கவையா அல்லவா என்று இப்போது கேட்கலாம்.

\end{document}
